\section{Physical tube volume and a normalized positive-source payment}
\label{sec:cv-tube-volume}
The total-variation estimate of Theorem~\ref{thm:ec-uniform-source}
does not by itself bound a physical-roof density.  This section proves
a comparison which does: an entire physical momentum disk can be
spread from each contributing critical tube, with its exact optical
weight bounded below as well as above.  Distinct tubes are disjoint
as physical sources.  Invariance then pays their complete volume
before the original-roof coarea is performed.  The result is uniform
in collision count and retains arbitrary Borel exact-record selectors.
It concerns the guarded critical-tube union already constructed in
Section~\ref{sec:pt-roof}; no coverage of its complement is assumed.

Keep $0<\varepsilon\le1/4$, the guarded regular critical roots
$\mathscr Z_{\varepsilon,m,R}$, and the positive constants
$s_\varepsilon,\rho_\varepsilon,C_{\rm dist}$ from
\eqref{eq:pt-momentum-radius} and \eqref{eq:pt-dist-constants}.
For $x\in\mathscr Z_{\varepsilon,m,R}$ let
\[
 \mathcal T_x=\{X_x(\zeta):|\zeta|<\rho_\varepsilon\},\qquad
 \mathcal T_{\varepsilon,m,R}
       =\bigcup_{x\in\mathscr Z_{\varepsilon,m,R}}\mathcal T_x,
\]
where $X_x(\zeta)$ is the actual physical initial collision state
of Lemma~\ref{lem:pt-momentum-disk}.  The whole tube is used here;
neither its initial nor its terminal section test is required to be
constant.  Put $u_x(\zeta)=|\partial_{\alpha_0}p_m(X_x(\zeta))|^{-1}$
and $u_x^0=u_x(0)$.  Throughout, the volume used for this auxiliary
comparison is the original unconditioned collision measure
$\nu=d\alpha\,dp/(4\pi)$.  The eventual source to be estimated still
uses exactly $\nu_R^*=c_*^{-1}\one_{Y_R^*}\nu$.

\begin{lemma}[Disjoint physical tubes and a two-sided optical volume]
\label{lem:cv-disjoint-volume}
Distinct members of $\mathscr Z_{\varepsilon,m,R}$ have disjoint
tubes as sets of physical initial states.  Each parametrization
$X_x$ is injective, and
\begin{equation}\label{eq:cv-distortion-volume}
 C_{\rm dist}^{-1}u_x^0\le u_x(\zeta)\le C_{\rm dist}u_x^0,
 \qquad
 \nu(\mathcal T_x)\ge
                   \frac{\rho_\varepsilon^2}{4C_{\rm dist}}u_x^0.
\end{equation}
Consequently, for any subfamily $\mathcal A$ whose entire tubes
are contained in a Borel set $E$ of the actual collision source,
\begin{equation}\label{eq:cv-coefficient-volume}
 \sum_{x\in\mathcal A}u_x^0
       \le\frac{4C_{\rm dist}}{\rho_\varepsilon^2}\nu(E).
\end{equation}
No whole-cylinder pressure or endpoint-kernel bound is used.
\end{lemma}
\begin{proof}
Every tube point is a regular first-hit trajectory with the same
strict lifted word as its center, by Theorem~\ref{thm:pt-physical}.
If two tubes contained the same physical initial state, uniqueness
of the actual successive first hits would give the same lifted word,
with the initial center normalized to zero.  Both centers would be
regular physical critical records of that word.
Corollary~\ref{cor:gv-critical} proves that such a record is the
unique minimum on the full product of the original closed disks.
Thus the centers coincide.  Lattice translates and angular seams
are identified in this argument, rather than counted as new centers.
Injectivity within a tube is also immediate from the fact that
$\zeta=(p_0,p_m)$ is an actual function of the physical initial
state; its local inverse has the specified values of both momenta.

The absolute logarithmic derivative in
Lemma~\ref{lem:pt-distortion} bounds the absolute difference
$|\log u_x(\zeta)-\log u_x^0|$ by
$D_*s_\varepsilon/\varepsilon=D_*\delta_*/2$ along the actual
physical endpoint segment.  This proves both inequalities in
\eqref{eq:cv-distortion-volume}, not just its upper inequality.
The exact inverse Jacobian from $(p_0,p_m)$ to $(\alpha_0,p_0)$
is $u_x$, as in Theorem~\ref{thm:rc-coarea}.  Hence
\[
 \nu(\mathcal T_x)
       =\frac1{4\pi}\int_{|\zeta|<\rho_\varepsilon}
                                  u_x(\zeta)\,d\zeta
       \ge\frac{\rho_\varepsilon^2}{4C_{\rm dist}}u_x^0.
\]
This identity has no section-conditioning factor and no additional
radius factor.  Sum it over finite subfamilies.  Their tubes are
disjoint physical sets contained in $E$, so their total mass is at
most $\nu(E)$.  Monotone exhaustion proves
\eqref{eq:cv-coefficient-volume} for the complete subfamily.
All inverse charts are smooth on their open domains and have Borel
images, so these are ordinary measurable source sets.
\end{proof}

\begin{lemma}[Every deep observation moves little across the whole tube]
\label{lem:cv-tube-observations}
Let $Q,Q'$ be any two physical records in one tube.  For
$d_i=\min(i,m-i)$ and $d'_i=\min(i,m-1-i)$ one has
\begin{align}
 |\Gamma_i(Q)-\Gamma_i(Q')|
       &\le L_\Gamma\sigma^{d'_i},&&0\le i<m,\notag\\
 |c_i(Q)-c_i(Q')|
       &\le L_c\sigma^{d_i},&&0\le i\le m,\label{eq:cv-observation-transport}\\
 d_\infty(x_i(Q),x_i(Q'))
       &\le L_\phi\sigma^{d_i},&&0\le i\le m,\notag
\end{align}
where
\begin{equation}\label{eq:cv-L}
 L_\Gamma=4C_\Delta s_\varepsilon,
 \qquad L_c=L_\phi=4C_c s_\varepsilon,
 \qquad\sigma=47/59.
\end{equation}
The clearance retains every third disk, including a disk visited
elsewhere in the word.  The distance in the last line is the
original circle--momentum distance.
\end{lemma}
\begin{proof}
Both endpoint pairs lie in the Euclidean ball of radius
$s_\varepsilon$ about the same critical center.  Their straight
endpoint segment lies in that ball, which consists of actual
physical trajectories by Theorem~\ref{thm:pt-physical}.
The length of this segment in the endpoint $\ell^\infty$ norm is
at most $2s_\varepsilon$.  Integrating
\eqref{eq:el-clearance} along it gives the first bound.  This
Lipschitz estimate already takes the infimum over the full allowed
set of centers, so switches of the minimizing center do not require
additional pieces or a multiplicity factor.

The incidence estimate follows in the same way from
\eqref{eq:el-incidence} and $s_i\le2\sigma^{d_i}$.
For clarity, the same derivative bound applies to the physical
tangential momentum.  For an outgoing velocity $v_i$ and the unit
tangent $t_i$, $p_i=v_i\cdot t_i$, so
\[
 |D_qp_i[\eta]|
       \le\|D_qv_i[\eta]\|_2+|D_q\alpha_i[\eta]|
       \le C_cs_i\|\eta\|_\infty.
\]
At the terminal contact use the incoming velocity; its tangential
component agrees with the outgoing one under specular reflection.
The adjacent-depth factor $\sigma^{-1}$ has already been included
in $C_c$.  Together with
$|D_q\alpha_i[\eta]|\le J_*s_i\|\eta\|_\infty$ and $J_*\le C_c$,
integration proves the last line.  Local angular lifts may be used
on the physical segment; projection to the circle can only decrease
the resulting distance.  Distance to the fixed set $\partial Y_R^*$
is one-Lipschitz in this same metric.
\end{proof}

For nonnegative depth widths $b_\Gamma(d),b_c(d),b_\phi(d)$ define
$\mathcal B_D(\boldsymbol b)$ on the original physical record by
\begin{align*}
 &\bigcup_{\substack{0\le i<m\\d'_i\ge D}}
                 \{\Gamma_i\le b_\Gamma(d'_i)\}\quad\cup\quad
   \bigcup_{\substack{0<i<m\\d_i\ge D}}
                 \{c_i\le b_c(d_i)\}\\
 &\hspace{24mm}\cup\quad
   \bigcup_{\substack{0<i<m\\d_i\ge D}}
      \{\operatorname{dist}_\infty(x_i,\partial Y_R^*)
                                             \le b_\phi(d_i)\}.
\end{align*}
These are source restrictions, not modifications of the return event.
Set
\begin{equation}\label{eq:cv-enlarged-widths}
 \widetilde b_\Gamma(d)=b_\Gamma(d)+L_\Gamma\sigma^d,
 \quad \widetilde b_c(d)=b_c(d)+L_c\sigma^d,
 \quad \widetilde b_\phi(d)=b_\phi(d)+L_\phi\sigma^d.
\end{equation}
Assume, for every $d\ge D$, that
\begin{equation}\label{eq:cv-width-regime}
 0\le\widetilde b_\Gamma(d)<3/200,\qquad
 0\le\widetilde b_c(d)\le1,\qquad
 0\le\widetilde b_\phi(d)\le1/200.
\end{equation}
Put, with nonnegative series interpreted by monotone limits,
\begin{equation}\label{eq:cv-unconditioned-budget}
 \mathcal V_D(\boldsymbol b)
  =\frac2{R_-}\sum_{d\ge D}\widetilde b_\Gamma(d)
       +2\sum_{d\ge D}\widetilde b_c(d)^2
       +\frac1\pi\sum_{d\ge D}\widetilde b_\phi(d).
\end{equation}

\begin{theorem}[A complete positive-source bound at the original roof]
\label{thm:cv-positive-roof}
For all $R,m,n,k$, $0<\varepsilon\le1/4$ and every Borel
$0\le h\le1$, the nonnegative measure
\[
 (L_{m,R})_*\left(
  \one_{\mathcal T_{\varepsilon,m,R}}
  \one_{\mathcal B_D(\boldsymbol b)}
  \one_{\Pi_{n,k,m,R}}h\,d\nu_R^*\right)
\]
is absolutely continuous with density $Q^{\rm bad}_{\varepsilon,D}$
satisfying, whenever \eqref{eq:cv-width-regime} holds,
\begin{equation}\label{eq:cv-positive-roof}
 \|Q^{\rm bad}_{\varepsilon,D}\|_{L^\infty(dt)}
   \le\frac{2U C_{\rm dist}^2}{c_*\rho_\varepsilon^2}
                                  \mathcal V_D(\boldsymbol b),
 \qquad U=320/9.
\end{equation}
The bound completes the full physical word sum before the
roof essential supremum or the radius supremum.  It uses neither
Input~\ref{input:v112-geometric-pressure} nor
Theorem~\ref{thm:mc-kernel}.
\end{theorem}
\begin{proof}
Select every tube which contains at least one point of
$\mathcal B_D(\boldsymbol b)$.  If a tube is selected, choose one
of its actual witnessed observations.  The same observation at
every other point of that tube satisfies its enlarged threshold
in \eqref{eq:cv-enlarged-widths}, by
Lemma~\ref{lem:cv-tube-observations}.  Thus the entire selected
tube lies in $\mathcal B_D(\widetilde{\boldsymbol b})$.
This reasoning does not require a witness common to different
tubes.  The family of centers is countable at fixed $R,m$, so
selection of a subfamily introduces no measurable-selection premise.
The enclosing event itself is a finite union of actual Borel
observations.

For the unconditioned physical measure, invariance and the exact
clearance flux, incidence distribution, and original section collar
give, respectively,
\[
 \nu\{\Gamma_i\le u\}\le u/R_-,\qquad
 \nu\{c_i\le u\}\le u^2,\qquad
 \nu\{\operatorname{dist}_\infty(x_i,\partial Y_R^*)\le u\}
                                                     \le u/(2\pi).
\]
The hypotheses for these three bounds are exactly
\eqref{eq:cv-width-regime}.  For clearance, closed inequality
boundaries are handled by continuity from above using slightly
larger widths still below $3/200$.  The exact incidence formula
and the closed-rectangle collar bound already include their
displayed upper endpoints; zero widths follow by continuity from
above.  At most two contacts or flights
have each specified depth.  The union bound therefore gives
\[
 \nu(\mathcal B_D(\widetilde{\boldsymbol b}))
                                      \le\mathcal V_D(\boldsymbol b).
\]
No independence of observations at different times is used.
Lemma~\ref{lem:cv-disjoint-volume} now proves
\[
 \sum_{\text{selected }x}u_x^0
    \le\frac{4C_{\rm dist}}{\rho_\varepsilon^2}
                                  \mathcal V_D(\boldsymbol b).
\]

On a selected tube use the exact polar coarea of
Theorem~\ref{thm:rc-coarea}.  At a fixed physical roof and angle,
there is at most one contributing radius in its momentum disk,
and its density is $u_x(\zeta)/(4\pi c_*h_x(\zeta))$.
Here $h_x\ge U^{-1}$ and $u_x\le C_{\rm dist}u_x^0$.
All source and exact-record selectors stay inside this density
until they are bounded above by one.  Summing the preceding
nonnegative coefficient bound and integrating an angle of length
$2\pi$ gives exactly \eqref{eq:cv-positive-roof}.
One first applies this argument to finitely many tubes and
nonnegative roof tests, then uses monotone convergence.  This proves
absolute continuity as well as its bound, without presupposing a
finite density.  The polar origins have zero collision area; no
measure supported on a singular paired boundary has been discarded.
The single factor $c_*^{-1}$ occurs only in this final original-source
coarea, not in the preceding unconditioned tube-volume comparison.
\end{proof}

\begin{corollary}[The exact-record exceptional source is paid on these tubes]
\label{cor:cv-ec-normalized}
Use the unchanged amplitudes of \eqref{eq:ec-amplitudes} and put
\[
 b_\Gamma(d)=A_\Gamma q_{\rm a}^{j(d)/2},\quad
 b_c(d)=A_cq_{\rm a}^{j(d)/2},\quad
 b_\phi(d)=A_\phi q_{\rm a}^{j(d)/4}.
\]
Then $\mathcal B_D(\boldsymbol b)=\mathcal E^{\rm obs}_{D,R,m}$.
For every $D\ge D_{\rm cv}:=20861442654$, independently of
$\varepsilon,R,m,n,k$, the width regime holds and
\begin{align}\label{eq:cv-explicit-budget}
 \mathcal V_D(\boldsymbol b)\le{}&
 \frac{2NA_\Gamma}{R_-}
          \frac{q_{\rm a}^{j(D)/2}}{1-q_{\rm a}^{1/2}}
 +\frac{2L_\Gamma}{R_-}\frac{\sigma^D}{1-\sigma}\notag\\
 &+4NA_c^2\frac{q_{\rm a}^{j(D)}}{1-q_{\rm a}}
 +4L_c^2\frac{\sigma^{2D}}{1-\sigma^2}\notag\\
 &+\frac{NA_\phi}{\pi}
          \frac{q_{\rm a}^{j(D)/4}}{1-q_{\rm a}^{1/4}}
 +\frac{L_\phi}{\pi}\frac{\sigma^D}{1-\sigma}.
\end{align}
For every fixed $0<\theta<1/2$, set $D_m=\lceil\theta m\rceil$.
At each fixed $\varepsilon>0$,
\begin{equation}\label{eq:cv-normalized-limit}
 \lim_{m\to\infty}\sup_{R,n,k,h}
             m^2\|Q^{\rm bad}_{\varepsilon,D_m}\|_\infty=0.
\end{equation}
In particular this pays, in the original normalized roof topology,
the entire exceptional bulk source of
Theorem~\ref{thm:ec-uniform-source} which belongs to
$\mathcal T_{\varepsilon,m,R}$.  It is a density theorem for a
specified full positive source, not a conversion of total variation
to essential height without a comparison argument.
\end{corollary}
\begin{proof}
Here is a rational certificate for the deliberately conservative
common lower depth.  The retained constants give
\[
 A_\Gamma\le12552,\quad A_c\le4,\quad A_\phi\le485,\quad
 L_\Gamma\le\frac{12756411}{9390862458880},\quad
 L_c=L_\phi\le\frac{8535}{284227072}.
\]
The last two bounds use $\varepsilon\le1/4$ and
$\delta_*=27/67443712$.  Writing $\kappa=3263680$, Bernoulli's
inequality gives
$(1+1/\kappa)^\kappa\ge2$, hence $q_{\rm a}^\kappa\le1/2$.
Also $2\cdot47^4<59^4$, so $\sigma^4<1/2$.
At $D_{\rm cv}=94(1+68\kappa)$ one has $j(D_{\rm cv})=68\kappa$,
and consequently
\[
 q_{\rm a}^{j(D_{\rm cv})/4}\le2^{-17},\qquad
 q_{\rm a}^{j(D_{\rm cv})/2}\le2^{-34},\qquad
 \sigma^{D_{\rm cv}}\le2^{-17}.
\]
Inserting the preceding rational amplitude bounds gives exactly
the three strict or weak inequalities in
\eqref{eq:cv-width-regime}; for example the collar bound is
$(485+8535/284227072)2^{-17}<1/200$.
All widths are nonincreasing in $D$, proving the stated common
regime without a floating-point logarithm or a depth chosen after
the collision count.
Use \eqref{eq:ec-depth-sum} for the $q_{\rm a}$ terms, ordinary
geometric summation for the $\sigma$ terms, and
$(a+b)^2\le2a^2+2b^2$ for the incidence term.  This proves all
six constants in \eqref{eq:cv-explicit-budget}.

For $D_m\ge N$, $j(D_m)\ge\theta m/N-2$.
Every term on its right is bounded by a constant times
$e^{-\lambda m}$, where one may take
\[
 \lambda=\theta\min\left\{
                  \frac{|\log q_{\rm a}|}{4N},|\log\sigma|\right\}>0.
\]
The prefactor in \eqref{eq:cv-positive-roof} is finite at fixed
$\varepsilon$.  Multiplication by $m^2$ proves
\eqref{eq:cv-normalized-limit}.  The bounds held at every finite
$m,D$ before this substitution.  Thus this depth classification
changes no microscopic roof window, physical clock, exact event,
or originally fixed outer parameter.  If $\varepsilon$ is a
positive quantity fixed at an outer stage $B$, the inner limit is
zero for each such $B$ before $B$ is sent to infinity.
\end{proof}

\begin{corollary}[The actual linear incidence multiplier and persistent source]
\label{cor:cv-incidence-persistent}
The preceding bounds remain valid with the source multiplied by
the original
$I_m=1+\sum_{i=1}^{m-1}q_0^{d_i}/c_i$, after multiplication of
their right sides by the explicit collision-uniform factor
\begin{equation}\label{eq:cv-I}
 H_\varepsilon=1+\frac4\varepsilon
                  \frac{q_0^{3/4}}{1-q_0^{3/4}},\qquad q_0=47/53.
\end{equation}
The same normalized payment applies to the Borel restriction of
the inherited persistent low-load interior source to
$\mathcal T_{\varepsilon,m,R}\cap\mathcal E^{\rm obs}_{D_m,R,m}$,
using its specific domination in Lemma~\ref{lem:v111-domination}.
\end{corollary}
\begin{proof}
Every tube point has the actual half-guard
$c_i\ge(\varepsilon/2)q_0^{d_i/4}$ by
Theorem~\ref{thm:pt-physical}.  Hence
\[
 I_m\le1+\frac2\varepsilon\sum_{i=1}^{m-1}q_0^{3d_i/4}
      \le1+\frac4\varepsilon
                    \sum_{d=1}^\infty q_0^{3d/4}=H_\varepsilon.
\]
This is a bound for the actual retained multiplier on these
physical tubes.  Positivity proves the weighted statement.
For the last assertion multiply
$0\le c_{\rm per}^{\varepsilon(B),\gamma_B}
       \le f^{\varepsilon(B)}\le A$
by the stated Borel restriction and original low-load selector.
Its actual source sum is dominated by the corresponding restricted
original roof measure.  Apply \eqref{eq:cv-normalized-limit}.
The tube width $\varepsilon$ may be fixed independently of the
inherited outer width $\varepsilon(B)$; if chosen from it, it must
be fixed before the collision limit.  No assertion is made about a
singular paired endpoint trace or a coefficient lacking this
specific domination.
\end{proof}

\begin{remark}[Exact scope and the still required full-source estimate]
\label{rem:cv-scope}
The new estimate pays actual bulk source at its common physical
roof through a proved lower optical volume and physical tube
disjointness.  Its dependence on the fixed tube radius is explicit:
$\rho_\varepsilon^{-2}$.  It establishes no uniform bound as that
radius tends to zero and therefore cannot be extended to every
regular trajectory merely by exhausting smaller guard classes.
Neither singular paired traces nor a full $I_m/|\nabla F|$
insertion is covered by \eqref{eq:cv-I}.
The original shallow and uncovered boundary sources, comparison
capacity, full physical multipliers, and the remaining spectral,
arithmetic and conditioning interfaces retain their exact
obligations.  In particular the original ordered limit
\eqref{eq:v112-remaining}, the raw pointwise local limit theorem,
and the four-coordinate inversion are not claimed closed here.
\end{remark}
